% Standalone problem document, generated from the adjacent README.md.
% Compile with XeLaTeX. Mathematical content is maintained in Markdown.
\documentclass[11pt,a4paper]{article}
\usepackage[fontsize=10.5pt]{scrextend}
\usepackage[margin=22mm,headheight=15pt,headsep=9mm,footskip=11mm]{geometry}
\usepackage{fontspec}
\setmainfont{texgyrepagella}[Extension=.otf,UprightFont=*-regular,BoldFont=*-bold,ItalicFont=*-italic,BoldItalicFont=*-bolditalic]
\setsansfont{texgyreheros}[Extension=.otf,UprightFont=*-regular,BoldFont=*-bold,ItalicFont=*-italic,BoldItalicFont=*-bolditalic]
\setmonofont{texgyrecursor}[Extension=.otf,UprightFont=*-regular,BoldFont=*-bold,ItalicFont=*-italic,BoldItalicFont=*-bolditalic,Scale=MatchLowercase]
\usepackage{amsmath,amssymb}
\usepackage{unicode-math}
\setmathfont{texgyrepagella-math.otf}
\usepackage{xcolor,microtype,enumitem,needspace,fancyhdr}
\usepackage{bookmark}
\definecolor{ink}{HTML}{17344B}
\definecolor{accent}{HTML}{197D80}
\definecolor{muted}{HTML}{596773}
\hypersetup{colorlinks=true,linkcolor=accent,urlcolor=accent,pdftitle={RA-10 - Constant-loss
nuclear-error transfer without matrix
ordering},pdfauthor={Open Problems in Numerical Linear Algebra}}
\urlstyle{same}
\setlength{\parindent}{0pt}
\setlength{\parskip}{5pt plus 1pt minus 1pt}
\linespread{1.04}
\setlength{\emergencystretch}{3em}
\widowpenalty=10000
\clubpenalty=10000
\displaywidowpenalty=10000
\allowdisplaybreaks[1]
\setlist{nosep,leftmargin=1.5em}
\providecommand{\tightlist}{\setlength{\itemsep}{0pt}\setlength{\parskip}{0pt}}
\providecommand{\pandocbounded}[1]{#1}
\pagestyle{fancy}
\fancyhf{}
\fancyhead[L]{\sffamily\footnotesize\color{muted}OPEN PROBLEMS IN NUMERICAL LINEAR ALGEBRA}
\fancyhead[R]{\sffamily\bfseries\footnotesize\color{accent}RA-10}
\fancyfoot[L]{\parbox[b]{0.9\textwidth}{\sffamily\fontsize{8}{10}\selectfont\color{muted}Editorial ratings; a bounded search cannot certify that a problem remains unsolved.\\Literature check: 2026-09-10}}
\fancyfoot[R]{\sffamily\footnotesize\color{muted}\thepage}
\renewcommand{\headrulewidth}{0.3pt}
\renewcommand{\headrule}{\hbox to\headwidth{\color{accent}\leaders\hrule height \headrulewidth\hfill}}
\makeatletter
\renewcommand{\section}{\Needspace{5\baselineskip}\@startsection{section}{1}{0pt}{12pt plus 2pt minus 2pt}{4pt}{\sffamily\bfseries\large\color{ink}}}
\renewcommand{\subsection}{\Needspace{4\baselineskip}\@startsection{subsection}{2}{0pt}{9pt plus 1pt minus 1pt}{3pt}{\sffamily\bfseries\normalsize\color{ink}}}
\makeatother
\setcounter{secnumdepth}{0}
\begin{document}
{\sffamily\small\color{muted}Randomized and low-rank approximation\par}
\vspace{3pt}
{\raggedright\hyphenpenalty=10000\exhyphenpenalty=10000\sffamily\bfseries\fontsize{21}{24}\selectfont\color{ink}Constant-loss
nuclear-error transfer without matrix ordering\par}
\vspace{7pt}
{\sffamily\small\textbf{Difficulty:} challenging\quad\textbf{Importance:} interesting
to the community\par}
{\sffamily\small\bfseries\color{accent}Status: Partially resolved\par}
\vspace{4pt}
{\color{accent}\hrule height 0.8pt}
\vspace{6pt}
\textbf{Rating rationale:} Challenging because a dimension-independent
loss must survive removal of matrix ordering; community impact is
transferring nuclear-error guarantees between matrix functions.\\
\textbf{Topic:} Low-rank approximation of matrix functions.

Does a universal constant \(C\ge1\) exist with the following property?
For every \(n\ge2\), \(1\le k<n\), real symmetric positive semidefinite
matrices \(A,\widehat A\in\mathbb R^{n\times n}\), \(\varepsilon\ge0\),
and continuous operator-monotone function \(f:[0,\infty)\to[0,\infty)\),

\[
\|A-\widehat A_k\|_*\le(1+\varepsilon)\|A-A_k\|_*
\]

implies

\[
\|f(A)-f(\widehat A)_k\|_*\le(1+C\varepsilon)\|f(A)-f(A)_k\|_*?
\]

Here \(\|M\|_*=\sum_i\sigma_i(M)\) is the nuclear norm. Operator
monotonicity means \(f(H)\succeq f(G)\) for every size and every real
symmetric \(H\succeq G\succeq0\); functions of matrices are defined
spectrally. No ordering between \(A\) and \(\widehat A\) is assumed.

For an ordered orthonormal eigendecomposition
\(X=\sum_{i=1}^n\lambda_iq_iq_i^T\) of a PSD matrix, set

\[
X_k=\sum_{i=1}^k\lambda_iq_iq_i^T,
\qquad f(X)_k=\sum_{i=1}^kf(\lambda_i)q_iq_i^T.
\]

The assertion must hold for every choice of ordered eigenvectors, using
the same choice in both truncations. \(C\) must be independent of all
inputs, including \(f\).

This is the nuclear-norm instance of the authors' explicit question
about a fixed loss in approximation accuracy. It asks whether the
ordering condition in existing transfer results can be removed at a
controlled price.

\newpage
\subsection{References}\label{references}

\begin{enumerate}
\def\labelenumi{\arabic{enumi}.}
\tightlist
\item
  D. Persson, R. A. Meyer, and C. Musco,
  \href{https://arxiv.org/html/2311.14023v2}{Algorithm-agnostic low-rank
  approximation of operator monotone matrix functions}, \emph{SIAM
  Journal on Matrix Analysis and Applications} 46 (2025), 1--21. §1.2,
  paragraph immediately preceding Table 1; §5, Example 5.3 and closing
  paragraph; Theorem 2.4 is the ordered predecessor. Equation (2)
  specifies truncation.
\item
  D. Persson, T. Chen, and C. Musco,
  \href{https://arxiv.org/abs/2502.01888}{Randomized block-Krylov
  subspace methods for low-rank approximation of matrix functions},
  \emph{Linear Algebra and its Applications} 741 (2026), 32--65,
  abstract and algorithmic scope.
\end{enumerate}

\subsection{Status check --- 2026-09-08}\label{status-check-2026-09-08}

Checked the source's latest listed arXiv v2 (2024-07-04), its 2025
journal record, and the exact counterexample and constant-loss
discussion. Searches combining the title, ``constant'',
``counterexamples'', ``nuclear'', ``funNyström'', and 2026 found no
general resolution. Example 5.3 rules out \(C=1\), while the source
expressly leaves a larger fixed constant possible. The 2026 Krylov work
addresses specific algorithms, not arbitrary PSD approximation pairs.
This entry is explicitly a specialization of the source's broader
question; the source does not number it as a separate conjecture. No
later explicit reaffirmation was located.

\subsection{Audit --- 2026-09-10}\label{audit-2026-09-10}

Rechecked \href{https://arxiv.org/html/2311.14023v2}{Theorem 2.4 and
Example 5.3}: ordered pairs satisfy \(C=1\), while unordered pairs can
violate that constant. A larger universal constant remains open.
Constant-loss nuclear-transfer searches found no resolution; the partial
tag refers only to the ordered subclass.
\end{document}
